% Compile with XeLaTeX, TeXLive 2013 or more recent
\documentclass{beamer}

% Base packages
\usepackage{fontspec}
\usepackage{xunicode}
\usepackage{xltxtra}

\usepackage{amsfonts}
\usepackage{amsmath}
\usepackage{longtable}
\usepackage{csquotes}
\usepackage{standalone}

\usepackage{graphicx}
\graphicspath{{./images/}}

\usepackage{listings}
\lstset{basicstyle=\footnotesize\ttfamily, breaklines=true, keepspaces=true }

% Setup Russian hyphenation
\usepackage{polyglossia}
\setdefaultlanguage[spelling=modern]{russian} % for polyglossia
\setotherlanguage{english} % for polyglossia
\defaultfontfeatures{Scale=MatchLowercase, Mapping=tex-text}

% Setup fonts
\newfontfamily\russianfont{CMU Serif}
\setromanfont{CMU Serif}
\setsansfont{CMU Sans Serif}
\setmonofont{CMU Typewriter Text}

% Be able to insert hyperlinks
\usepackage{hyperref}
\hypersetup{colorlinks=true, linkcolor=black, filecolor=black, citecolor=black, urlcolor=blue , pdfauthor=Evgeny Yulyugin <yulyugin@gmail.com>, pdftitle=Параллельное программирование}
% \usepackage{url}

% Misc optional packages
\usepackage{underscore}
\usepackage{amsthm}

% A new command to mark not done places
\newcommand{\todo}[1][]{{\color{red}TODO\ #1}}

\newcommand{\abbr}{\textit{англ.}\ }

\subtitle{Курс «Параллельное программирование»}
\subject{Lecture}
\author[Евгений Юлюгин]{Евгений Юлюгин \\ \small{\href{mailto:yulyugin@gmail.com}{yulyugin@gmail.com}}}
\date{\today}
\pgfdeclareimage[height=0.5cm]{mipt-logo}{../common-images/mipt.png}
\logo{\pgfuseimage{mipt-logo}}

\typeout{Copyright 2014 Evgeny Yulyugin}

\usetheme{Berlin}
\setbeamertemplate{navigation symbols}{}%remove navigation symbols

\title{Основы MPI}

\begin{document}

\begin{frame}
\titlepage
\end{frame}

\section{Обзор}

\begin{frame}
\tableofcontents
\end{frame} 

\section{История MPI}

\begin{frame}{История MPI}

\begin{itemize}
	\item MPI расшифровывается как интерфейс передачи сообщений (\abbr Massage Passing Interface).
	\item Программный интерфейс, позволяющий обмениваться сообщениями между процессами выполняющими одну задачу.
	\item Разработан Уильямом Гроуппом, Эвином Ласком и другими.
	\item Первая версия разрабатывалась в 1993-1994 году.
	\item MPI версии 1 вышла в 1994 году.
	\item MPI 1.1 опубликован в 12 июня 1995 года. Поддерживается большинством современных реализаций MPI. Первая раелизация появилась в 2002 году.
	\item Существуют реализации для языков Fortran, Java, C и C++.
\end{itemize}

\end{frame}

\section{Базовые функции MPI}

\defverbatim[colored]\initfinalize{
\begin{lstlisting}[language=C++,basicstyle=\ttfamily,keywordstyle=\color{blue}]
int MPI_Init(int *pargc, char ***pargv);
int MPI_Finalize(void);
\end{lstlisting}
}

\begin{frame}{Инициализация и завершение процессов}

Определены в заголовочном файле mpi.h

\vfill

\initfinalize

\end{frame}

\defverbatim[colored]\size{
\begin{lstlisting}[language=C++,basicstyle=\ttfamily,keywordstyle=\color{blue}]
MPI_Comm_size(MPI_Comm comm, int *size);
\end{lstlisting}
}

\defverbatim[colored]\rank{
\begin{lstlisting}[language=C++,basicstyle=\ttfamily,keywordstyle=\color{blue}]
MPI_Comm_rank(MPI_Comm comm, int *rank);
\end{lstlisting}
}

\begin{frame}{MPI communicator}

\begin{itemize}
	\item MPI_COMM_WORLD is the initially defined universe intracommunicator for all processes to conduct various communications once MPI_INIT has been called,
	\item MPI_COMM_SELF,
	\item User defined communicators.
\end{itemize}

\end{frame}

\begin{frame}{Self-identification}

\size

Returns the size of the group associated with communicator.

\vfill

\rank

Determines the rank of the calling process in the communicator.

\end{frame}

\defverbatim[colored]\sendrecv{
\begin{lstlisting}[language=C++,basicstyle=\ttfamily,keywordstyle=\color{blue}]
int MPI_Send(void *buf, int count, MPI_Datatype datatype, int dest, int tag, MPI_Comm comm);
int MPI_Recv(void *buf, int count, MPI_Datatype datatype, int source, int tag, MPI_Comm comm, MPI_Status status);
\end{lstlisting}
}

\begin{frame}{Send/Receive Overview}

\sendrecv

\end{frame}

\begin{frame}{MPI_Datatype}

\begin{itemize}
	\item MPI_CHAR,
	\item MPI_INT,
	\item MPI_UNSIGNED,
	\item MPI_DOUBLE,
	\item etc.,
	\item Custom data types.
\end{itemize}

\end{frame}

\begin{frame}{MPI_Status \& tag}

MPI_Status:

\begin{itemize}
	\item Содержит дополнительную информацию о полученном сообщении.
	\item MPI_STATUS_IGNORE --- специальное значение, которое уменьшает количество используемой памяти в том случае, если вы не собираетесь рассматривать это поле.
\end{itemize}

\vfill

tag:

\begin{itemize}
	\item Число для идентификации сообщения,
	\item MPI_ANY_TAG.
\end{itemize}

\end{frame}

\begin{frame}{Блокирующие и не блокирующие операции}

MPI_Send/Recv --- блокирующиеся функции.

\begin{figure}[htp]
	\centering
	\includegraphics[width=0.4\textwidth]{blocked}
\end{figure}

MPI_Isend/Irecv – не блокирующиеся варианты MPI_Send/Recv.

\end{frame}

\defverbatim[colored]\wtime{
\begin{lstlisting}[language=C++,basicstyle=\ttfamily,keywordstyle=\color{blue}]
double MPI_Wtime()
\end{lstlisting}
}

\defverbatim[colored]\wtick{
\begin{lstlisting}[language=C++,basicstyle=\ttfamily,keywordstyle=\color{blue}]
double MPI_Wtick()
\end{lstlisting}
}

\defverbatim[colored]\barrier{
\begin{lstlisting}[language=C++,basicstyle=\ttfamily,keywordstyle=\color{blue}]
int MPI_Barrier(MPI_Comm comm)
\end{lstlisting}
}

\begin{frame}{Вычисление времени}

\wtime

Возвращает время в секундах, прошедшее с  некоторого момента в прошлом (точки отсчета). Гарантируется, что эта точка отсчета не будет изменена в течение жизни процесса.

\wtick

Возвращает разрешение таймера (минимальное значение кванта времени). 

\barrier

Блокирует работу вызвавшего ее процесса до тех пор, пока все другие процессы группы также не вызовут эту функцию.

\end{frame}

\section{Компиляция и запуск}

\defverbatim[colored]\mpirun{
\begin{lstlisting}[language=bash,basicstyle=\ttfamily,keywordstyle=\color{blue}]
$ mpirun –np <n_processes> <programm> <args>
\end{lstlisting}
}

\begin{frame}{Компиляция и запуск}

Для компиляции используется \texttt{mpicc} вместо \texttt{gcc}.

\vfill

Запуск производится командой вида

\mpirun

\end{frame}

\section{Конец}
% The final "thank you" frame 

\begin{frame}{Задания}

\begin{itemize}
    \item MPI "Hello, World!",
    \item MPI ring,
    \item MPI ping-pong.
\end{itemize}

\end{frame}

\begin{frame}{На следующей лекции}

\begin{itemize}
    \item Классификация архитектур вычислительных систем,
    \item Состояние гонки,
    \item Примитивы синхронизации.
\end{itemize}

\end{frame}

\begin{frame}

{\huge{Спасибо за внимание!}\par}

\vfill

\tiny{\textit{Замечание}: все торговые марки и логотипы, использованные в данном материале, являются собственностью их владельцев. Представленная здесь точка зрения отражает личное мнение автора, не выступающего от лица какой-либо организации.}

\end{frame}

\end{document}
